\section{Bundle tests of bounded degree}
\label{sec:bundle-tests}

We now construct the finite topological tests used to exclude polynomial
witnesses of proper infiniteness.  The matrix size $m$ and the degree bound
$R$ determine the bundles and their base, while the number of maps remains
fixed at $24$.
The use of characteristic classes to obstruct comparison follows the
methods of Villadsen \cite[Section~3]{Villadsen1999} and R{\o}rdam
\cite[Sections~3--4]{Rordam2003}. Our test retains enough of a Chern
class to exclude a specified trivial subbundle, even though the maps
between successive levels have the normalized column and row bounds.

\begin{proposition}\label{prop:bundle-test}
For every pair of integers $m,R\geq 1$, set
\begin{equation}\label{eq:test-ranks}
 d_0=64m4^R,\qquad d_j=4^j d_0\quad(0\leq j\leq R).
\end{equation}
There are a compact smooth manifold $Y$, Hermitian complex vector bundles
$E_j\to Y$ for all $j\geq0$, and continuous bundle maps
$X_i^{(j)}:E_j\to E_{j+1}$, for $1\leq i\leq24$ and $j\geq0$, such that
$E_0$ is trivial of rank $d_0$, the rank of $E_j$ is $d_j$ for
$0\leq j\leq R$, and
\begin{equation}\label{bundle:scaled-identities}
 \sum_{i=1}^{24}(X_i^{(j)})^*X_i^{(j)}=I_{E_j},
 \qquad
 \sum_{i=1}^{24}X_i^{(j)}(X_i^{(j)})^*
 \leq\frac23 I_{E_{j+1}}.
\end{equation}
Moreover,
\begin{equation}\label{eq:chern-obstruction}
 \begin{gathered}
 \left(\bigoplus_{j=0}^R E_j\right)^{\oplus m}\\
 \text{contains no trivial subbundle of rank }2md_0.
 \end{gathered}
\end{equation}
\end{proposition}

\begin{proof}
Fix the matrix size $m$ and degree bound $R$.
The bundles must satisfy two requirements.  Their sum through level $R$
must retain a Chern class after allowing the trivial directions coming
from $E_0$.  At the same time, successive bundles must admit a fixed
number of maps that expand every source subspace.  We reduce each old
line-bundle multiplicity by a factor of four at every step, so the total
trivial part stays small.  Each new projective-space factor is then sized
to retain the Chern class contributed by every later occurrence of its
line bundle.  The initial rank leaves enough room in the reduced blocks
to expand small subspaces; a large new block handles the others.

\paragraph{The bundles and their base.}
For $1\leq i\leq R$, define the fresh multiplicities
\[
 b_i=\frac{15}{4}d_{i-1}.
\]
Let the multiplicity of the line bundle $L_i$ in $E_j$ be
\begin{equation}\label{bundle:multiplicities}
 \begin{split}
 a_{0,j}&=d_0 4^{-j}=64m4^{R-j},\\
 a_{i,j}&=b_i4^{-(j-i)}
          =960m4^{R+2i-j-2}\qquad(1\leq i\leq j\leq R).
 \end{split}
\end{equation}
All these numbers are positive integers.  Put
\begin{equation}\label{eq:test-base}
 K_i=m\sum_{j=i}^R a_{i,j}
     =320m^2 4^{2i-2}\bigl(4^{R-i+1}-1\bigr),
 \qquad
 Y=\prod_{i=1}^R\mathbb{CP}^{K_i}.
\end{equation}
Thus every $K_i$ is a positive integer.  Take $L_0$ to be the trivial line
bundle on $Y$, and for $i\geq1$ let $L_i$ be the pullback of the hyperplane
line bundle from the $i$th factor.  Choose Hermitian metrics and define
\[
 E_j=\bigoplus_{i=0}^j L_i^{\oplus a_{i,j}}
 \qquad(0\leq j\leq R).
\]
In passing from $E_j$ to $E_{j+1}$, each old multiplicity is divided by
four and a fresh block of rank $b_{j+1}=(15/4)d_j$ is added.
Consequently the ranks satisfy
\[
 \operatorname{rank}E_{j+1}
 =\frac14d_j+\frac{15}{4}d_j=4d_j,
\]
which proves the rank assertion by induction.  Formula~\eqref{bundle:multiplicities}
also shows that every multiplicity $D$ to which division by four is
applied is divisible by four and satisfies $D\geq256m$.

\paragraph{The retained Chern class.}
The choice of $K_i$ records exactly the total multiplicity of $L_i$ in
the bundle to be tested.  Set
\[
 G=\left(\bigoplus_{j=0}^R E_j\right)^{\oplus m},\qquad
 h_0=m\sum_{j=0}^R a_{0,j},\qquad p=\sum_{i=1}^R K_i.
\]
The definitions of the multiplicities and $K_i$ give
\begin{equation}\label{bundle:chern-decomposition}
 G\cong L_0^{\oplus h_0}\oplus
             \bigoplus_{i=1}^R L_i^{\oplus K_i},\qquad
 h_0=\frac43md_0\bigl(1-4^{-(R+1)}\bigr)<2md_0.
\end{equation}
Let $z_i=c_1(L_i)$.  The integral cohomology ring is
\[
 H^*(Y;\mathbb Z)
 \cong\mathbb Z[z_1,\ldots,z_R]/
       (z_1^{K_1+1},\ldots,z_R^{K_R+1}),\qquad \deg z_i=2.
\]
In particular $\prod_i z_i^{K_i}$ is nonzero.  The Whitney sum formula
and the normalization for line bundles give
\[
 c_p(G)=\prod_{i=1}^R z_i^{K_i}\ne0;
\]
we use the standard Chern-class properties in
\cite[Theorem~3.2]{HatcherVB}.
If $G$ had a trivial subbundle of rank $k>h_0$, its Hermitian complement
would have rank $h_0+p-k<p$.  The Whitney sum formula would identify
$c_p(G)$ with the $p$th Chern class of that complement, which vanishes
above its rank.  This contradiction, together with
\eqref{bundle:chern-decomposition}, proves \eqref{eq:chern-obstruction}.
The argument applies to any continuous subbundle, without a condition
on its position relative to the displayed line summands.

\paragraph{The bases for the expanding maps.}
It remains to construct the maps between these bundles.  For the step
$j\to j+1$, with $0\leq j<R$, the source and target use only
$L_0,\ldots,L_{j+1}$ and are pulled back from
\[
 Y_{j+1}=\prod_{i=1}^{j+1}\mathbb{CP}^{K_i}.
\]
The dimensions $K_i$ here are those in \eqref{eq:test-base}, including
all later occurrences of $L_i$ through level $R$.  The complex dimension
$\eta_j$ of this smaller base satisfies
\begin{equation}\label{bundle:prefix-dimension}
 \begin{split}
 \eta_j=\sum_{i=1}^{j+1}K_i
 &\leq\frac43m\sum_{i=1}^{j+1}b_i\\
 &=\frac53md_0(4^{j+1}-1)
 <\frac{20}{3}md_j.
 \end{split}
\end{equation}
Here the first inequality uses the geometric series with ratio $1/4$.
We construct the maps for this step on $Y_{j+1}$, and then pull them back
to $Y$.  The estimate by a multiple of $md_j$ is what will allow the
incidence argument to handle all subspaces of dimension greater than
$m$, even though the full base $Y$ may be much larger.

\paragraph{Sixteen maps for small subspaces.}
Decompose $E_{j+1}$ as the sum of its replacement blocks and its fresh
block.  On a source block $L_i^{\oplus D}$, the corresponding replacement
block is $L_i^{\oplus(D/4)}$.  We seek $16$ scalar matrices between the
multiplicity spaces and tensor them with the identity of $L_i$.
For $1\leq s\leq m$, set $r=D/4$ and $t=2s$.  Since $D\geq64m$,
\[
 t<r,\qquad r-2s\geq D/8.
\]
The incidence expression of Lemma~\ref{lem:incidence}, with $q=16$ and
base a point, is strictly positive:
\begin{equation}\label{bundle:small-incidence}
 (16s-2s)(D/4-2s)-s(D-s)
 =\frac52sD-27s^2
 \geq133ms>0.
\end{equation}
Applying Lemma~\ref{lem:incidence} simultaneously for $s=1,\ldots,m$ gives matrices
whose images of every such $s$-plane span a space of dimension at least
$2s$.  In fact, avoiding dimension at most $2s$ gives a slightly stronger
conclusion.  Make these choices for each block and assemble them into
$16$ block diagonal maps $T_1,\ldots,T_{16}$ to the replacement part,
with zero component in the fresh block.

We verify that the same $16$ maps expand every subspace of the whole
source of dimension at most $m$.  Work in one fiber and write the
ordered source and replacement blocks as $V_k$ and $F_k$.  Let
$V_{\leq k}$ and $F_{\leq k}$ denote their initial sums.  Given a subspace
$S$ of dimension at most $m$, set
\[
 S_k=S\cap V_{\leq k},\qquad
 W=\sum_{i=1}^{16}T_iS,\qquad
 W_k=W\cap F_{\leq k}.
\]
Use the zero spaces before the first flag terms.  Projection onto the
$k$th block identifies $S_k/S_{k-1}$ and $W_k/W_{k-1}$ with subspaces
of $V_k$ and $F_k$, respectively.  Block diagonality gives
$T_i(S_k)\subseteq W_k$ and hence
\[
 \sum_{i=1}^{16}T_{i,k}(S_k/S_{k-1})
 \subseteq W_k/W_{k-1},
\]
where $T_{i,k}$ is the map on the $k$th block.  Each source quotient has
dimension at most $m$, so the block estimates yield
\begin{equation}\label{bundle:flag-expansion}
 \dim W
 =\sum_k\dim(W_k/W_{k-1})
 \geq2\sum_k\dim(S_k/S_{k-1})
 =2\dim S.
\end{equation}
This proves the required estimate even when $S$ has correlated
components in several blocks.

\paragraph{Eight maps for large subspaces.}
Let $d=d_j$ and let the target now be the fresh block of rank
$r=(15/4)d$.  For every integer $m<s\leq d$, take $q=8$ and $t=2s$.
The rank endpoint is valid because $2s\leq2d<r$.  The incidence
expression satisfies
\begin{equation}\label{bundle:large-incidence}
 \begin{split}
 (8s-2s)(15d/4-2s)-s(d-s)
 &=\frac{43}{2}sd-11s^2\\
 &\geq\frac{21}{2}sd
 >\frac{20}{3}md
 >\eta_j.
 \end{split}
\end{equation}
The first inequality follows from $s\leq d$.  Lemma~\ref{lem:incidence} on $Y_{j+1}$, applied to the finite
list $s=m+1,\ldots,d$, therefore supplies eight smooth maps to the fresh
block whose image spans have dimension at least $2s$ for all these
subspaces.  Regard these as maps to $E_{j+1}$ with zero replacement
components.

Together, the two lists give $24$ maps $T_i:E_j\to E_{j+1}$ satisfying
\[
 \dim\sum_{i=1}^{24}T_iS\geq2\dim S
\]
for every fiber subspace $S$.  The estimate is supplied by the first
list when $\dim S\leq m$ and by the second when $\dim S>m$; adding maps
cannot decrease an image span.  Pull back to $Y$ and apply
Lemma~\ref{lem:scaling}.  This gives continuous maps $X_i^{(j)}$
satisfying \eqref{bundle:scaled-identities} for every $0\leq j<R$.
The scaling may mix the displayed line summands, but each resulting
map still has source $E_j$ and target $E_{j+1}$.

\paragraph{Continuation beyond the test.}
For $j\geq R$, define $E_{j+1}=E_j\oplus E_j$, and let
\[
 X_1^{(j)}=2^{-1/2}\iota_1,\qquad
 X_2^{(j)}=2^{-1/2}\iota_2,\qquad
 X_i^{(j)}=0\quad(3\leq i\leq24),
\]
where $\iota_1,\iota_2$ are the two inclusions.  Their column sum is
$I_{E_j}$ and their row sum is $\tfrac12I_{E_{j+1}}$, so
\eqref{bundle:scaled-identities} holds for these steps as well.
This continuation leaves the retained bundle in
\eqref{eq:chern-obstruction} unchanged and completes the construction.
\end{proof}
